\chapter{Supersimmetria \texorpdfstring{$\NSC = 1$}{N = 1} globale in \texorpdfstring{$D = 4$}{D = 4}}
In supersimmetria l'insieme delle simmetrie incluse in teoria quantistica dei campi è estesa dalle trasformazioni di simmetria interne e di Poincaré all'inclusione delle supercariche spinoriali $Q_\alpha^i$, dove $\alpha$ è un indice spinoriale e $i = 1,\dots,\fan{N}$ è un indice che etichetta le differenti supercariche. Si assume che $Q_\alpha^i$ siano spinori di Majorana a quattro componenti, anche se una formulazione equivalente tramite spinori di Weyl a due componenti è pure possibile. Ci si concentra per ora sul caso più semplice, cioè per $\NSC = 1$ dove c'è solo una singola supercarica $Q_\alpha$. Questo caso è chiamato \emph{$\NSC = 1$ SUSY} o \emph{simple SUSY}.

In supersimmetria $\NSC = 1$ globale i generatori di Poincaré e $Q_\alpha$ si incontrano in una nuova struttura algebrica, quella di una \emph{superalgebra}. Una superalgebra contiene due classi di elementi, pari e dispari. Da un punto di vista fisico, gli elementi delle due classi sono chiamati rispettivamente bosonici ($B$) e fermionici ($F$). Le relazioni di struttura includono sia commutatori che anti-commutatori con le regole generali $\comm{B}{B} = B$, $\comm{B}{F} = B$ e $\acomm{F}{F} = B$. Le cariche bosoniche definiscono un'algebra di Lie. In supersimmetria la superalgebra delle cariche bosoniche $M_{[\mu\nu]}$ e $P_\mu$ è l'algebra di Lie del gruppo di Poincaré\footnote{Le cariche $M_{[\mu\nu]}$ sono definite come $M_{[\mu\nu]} = \int d^{D - 1}\v{x}\,M^0_{[\mu\nu]}$.}, mentre le nuove relazioni di struttura che includono $Q_\alpha$ sono 
\begin{equation}
\label{eq: superalgebra per N = 1}
    \begin{cases}
        \acomm{Q_\alpha}{\t{Q}^\beta} = -\dfrac{1}{2}(\gamma_\mu)_\alpha^\beta P^\mu \\[2ex] 
        \comm{M_{[\mu\nu]}}{Q_\alpha} = -\dfrac{1}{2}(\gamma_{\mu\nu})_\alpha^\beta Q_\beta \\[2ex]
        \comm{P_\mu}{Q_\alpha} = 0
    \end{cases}.
\end{equation}
Molte delle teorie supersimmetriche, ma non tutte, sono invarianti sotto una trasformazione chirale $\group{U}{1}$ chiamata \emph{R-simmetria}. Si indica il generatore con $T_R$. Questo agisce su $Q_\alpha$ come 
\begin{equation}
    \comm{T_R}{Q_\alpha} = -i(\gamma_*)^\beta_\alpha Q_\beta,
\end{equation}
ma questo generatore $T_R$ non è richiesto. Altre simmetrie interne, che commutano con $W_\alpha$, frequentemente chiamate \emph{cariche esterne}, possono essere incluse.

I parametri delle trasformazioni di supersimmetria globale sono spinori di Majorana costanti anti-commutanti $\epsilon_\alpha$. In supergravità la supersimmetria ha il gauge fissato, necessariamente con i generatori di Poincaré, dato che sono inclusi nella superalgebra \eqref{eq: superalgebra per N = 1}. Questo implica che la gravità è inclusa, per cui gli spinori diventano delle funzioni arbitrarie $\epsilon_\alpha(x)$ su una varietà spaziotemporale curva. 

\section{Teorie supersimmetriche minimali}
Le teorie supersimmetriche contengono sia bosoni che fermioni, che sono gli stati di base di una rappresentazione di particella della superalgebra della supersimmetria \eqref{eq: superalgebra per N = 1}. Senza essere troppo tecnici, gli stati delle particelle com momento $\v{p}$ ed energia $E(\v{p}) = \sqrt{\v{p}^2 + m_{B,F}^2}$ sono indicati con $\ket{\v{p},B}$ e $\ket{\v{p},F}$, dove le etichette $B$ ed $F$ includono l'elicità delle particelle. Le trasformazioni di supersimmetria connettono questi stati. Dato che lo spinore $Q_\alpha$ porta con se momento angolare $1/2$, questo trasforma bosoni in fermioni, e viceversa. Quindi, $Q_\alpha\ket{\v{p},B} = \ket{\v{p},F}$ e $Q_\alpha\ket{\v{p},F} \propto \ket{\v{p},B}$. Dato che $\comm{P^\mu}{Q_\alpha} = 0$, gli stati trasformati hanno lo stesso momento ed energia, per cui la stessa, quindi $m_B^2 = m_F^2$. Si può dimostrare che una rappresentazione dell'algebra contiene lo stesso numero di bosoni e fermioni.

Le rappresentazioni più semplici dell'algebra che portano alle teorie di campo supersimmetriche più semplici, dette per ciò \emph{minimali}, sono:
\begin{enumerate}
    \item il \emph{multipletto chirale}, il quale consiste in un fermione auto-aggiunto di spin $1/2$ descritto da un campo di Majorana $\chi(x)$ e un bosone complesso di spin $0$ descritto da un campo scalare $Z(x)$. Alternativamente, il campo $\chi(x)$ può essere sostituito con lo spinore di Weyl $P_L\chi(x)$ e/o $Z(x)$ tramite la combinazione $Z(x) = [A(x) + B(x)]/\sqrt{2}$, dove $A(x)$ e $B(x)$ sono un campo scalare reale e un campo pseudo-scalare rispettivamente;
    \item il \emph{multipletto di gauge}, il quale consiste in una particelle massless di spin $1$ descritta da un campo vettoriale di gauge $A_\mu(x)$ e la sua controparte fermionica di spin $1/2$, il \emph{gaugino}, descritto da uno spinore di Majorana $\lambda(x)$ (o il corrispondente campo di Weyl $P_L\lambda(x)$).
\end{enumerate}

\subsection{Teorie supersimmetriche per il multipletto chirale}
Come detto più sopra, i campi fisici del multipletto chirale sono il campo scalare complesso $Z(x)$ e lo spinore di Majorana $\chi(x)$. L'introduzione di un campo scalare complesso ausiliario $F(x)$ semplifica la struttura in più modi. Le equazioni di campo del campo $F$ sono algebriche, per cui $F$ può essere eliminato dal sistema in un secondo momento. L'insieme dei campi $Z$, $P_L\chi$ ed $F$ costituisce il multipletto chirale, mentre i loro coniugati $\b{Z}$, $P_R\chi$ e $\b{F}$ formano il multipletto anti-chirale. 
\begin{notz}
    Con $\b{Z}$ e $\b{F}$ si indicano i complessi coniugati $Z^*$ e $F^*$, mentre con $\b{\epsilon}$ si indica effettivamente il coniugato di Majorana di $\epsilon$.
\end{notz}
Si vuole ora presentare le regole di trasformazioni di questi multipletti, discutere le azioni invarianti e studiare la superalgebra \eqref{eq: superalgebra per N = 1}. Il parametro spinoriale $\epsilon$ è uno spinore di Majorana, le cui componenti anti-commutano tra di loro e con le componenti di $\chi$ e $\b{\chi}$. Le regole di trasformazione per il multipletto chirale sono
\begin{equation}
\label{eq: trasformazioni multipletto chirale}
    \begin{cases}
        \delta Z = \dfrac{1}{\sqrt{2}}\b{\epsilon}P_L\chi \\[2ex]
        \delta P_L\chi = \dfrac{1}{\sqrt{2}}P_L(\sl{\p}Z + F)\epsilon \\[2ex]
        \delta F = \dfrac{1}{\sqrt{2}}\b{\epsilon}\sl{\p}P_L\chi
    \end{cases},
\end{equation}
mentre quelle del multipletto anti-chirale sono
\begin{equation}
\label{eq: trasformazioni multipletto anti-chirale}
    \begin{cases}
        \delta \b{Z} = \dfrac{1}{\sqrt{2}}\b{\epsilon}P_R\chi \\[2ex]
        \delta P_R\chi = \dfrac{1}{\sqrt{2}}P_R(\sl{\p}\b{Z} + \b{F})\epsilon \\[2ex]
        \delta \b{F} = \dfrac{1}{\sqrt{2}}\b{\epsilon}\sl{\p}P_R\chi
    \end{cases}.
\end{equation}
Ci sono due azioni fondamentali, ciascuna invariante in maniera indipendente sotto le regole di trasformazioni appena elencate. La prima è l'azione cinetica libera
\begin{equation}
    S_K = -\int d^4x\,(\p^\mu\b{Z}\p_\mu Z + \b{\chi}\sl{\p}P_L\chi - \b{F}F),
\end{equation}
dove si è incluso il termine spinoriale in forma chirale. L'interazione è determinata da una funzione arbitraria olomorfa $W(Z)$ detta \emph{superpotenziale}. Con quest'ultima, si definisce l'azione 
\begin{equation}
    S_F = \int d^4x\,\bigg[FW'(Z) - \frac{1}{2}\b{\chi}P_LW''(Z)\chi\bigg].
\end{equation}
Si noti che l'azione $S_F$ contiene solamente i campi del multipletto chirale e nessuno di quello anti-chirale; pertanto, $S_F$ non è hermitiana e si deve quindi considerare anche l'azione coniugata $S_{\b{F}} \equiv S_F^\dag$. Dunque, l'azione completa del multipletto chirale è 
\begin{equation}
    S = S_K + S_F + S_{\b{F}}.
\end{equation}
Si discute adesso in maniera informale la dimostrazione del fatto che le azioni $S_K$ ed $S_F$ sono invarianti sotto le trasformazioni di supersimmetria \eqref{eq: trasformazioni multipletto chirale} e \eqref{eq: trasformazioni multipletto anti-chirale}. Per quanto riguarda $S_K$, la dimostrazione viene semplificata da un'osservazione valida in qualsiasi rappresentazione dell'algebra di Clifford, ma che risulta più evidente nella rappresentazione di Weyl. Le proiezioni $P_L\epsilon$ e $\b{\epsilon}P_L$ coinvolgono la stessa metà delle componenti dello spinore di Majorana $\epsilon$, mentre $P_R\epsilon$ e $\b{\epsilon}P_R$ coinvolgono le componenti coniugate. La variazione totale si può scrivere come $\delta S = \delta_{P_L\epsilon}S + \delta_{P_R\epsilon}S$ separando temporaneamente le due proiezioni chirali di $\epsilon$ nelle regole di trasformazioni. Dato che $S_K$ è hermitiana, è sufficiente calcolare $\delta_{P_L\epsilon}S_K$, dato che $\delta_{P_R\epsilon}S$ è il suo aggiunto. Nel calcolo si assume temporaneamente che $\epsilon(x)$ sia una funzione arbitraria nello spaziotempo di Minkowski in modo da ottenere in maniera semplice la corrente di Noether per la supersimmetria. Serve anche $\delta\b{\chi}P_R \equiv \delta \chi^TCP_R$, che usando
\begin{equation}
    (\delta P_R\chi)^TC = \delta \chi^TP_R^TC = \delta \chi^TCC^{-1}P_R^TC,
\end{equation}
si può dimostrare valere 
\begin{equation}
    \delta \b{\chi}P_R = \frac{1}{\sqrt{2}}\b{\epsilon}(-\sl{\p}\b{Z} + \b{F})P_R.
\end{equation}
Avendo impostato la discussione, si calcola 
\begin{equation}
\label{eq: variazione chirale left azione cinetica}
    \begin{split}
        \delta_{P_L\epsilon}S_K &= \frac{1}{\sqrt{2}}\int d^4x\,[-\p^\mu \b{Z}\p_\mu(\epsilon P_L\chi) + \b{\epsilon}P_L(\sl{\p}\b{Z})\sl{\p}P_L\chi \\
        &- \b{\chi}P_R\sl{\p}(P_LF\epsilon) + F\b{\epsilon}P_L\sl{\p}\chi].
    \end{split}
\end{equation}
I termini in $\b{Z}\chi$ e $F\chi$ sono indipendenti e devono annullarsi separatamente nel caso di una simmetria, se $\epsilon$ è costante. Dopo un Majorana flip sull'ultimo termine, si trova che i termini $F\chi$ si combinano in 
\begin{equation}
    -\frac{1}{\sqrt{2}}\int d^4x\,\p_\mu(\b{\chi}\gamma^\mu P_LF\epsilon),
\end{equation}
che si annulla. I termini in $\b{Z}\chi$ si possono manipolare sostituendo le relazioni
\begin{equation}
    \p_\mu(\b{\epsilon}P_L\chi) = (\p_\mu\b{\epsilon})P_L\chi + \b{\epsilon}P_L\p_\mu\chi,    
\end{equation}
\begin{equation}
    \b{\epsilon}P_L\gamma^\mu\gamma^\nu(\p_\mu\b{Z})\p_\nu\chi = \b{\epsilon}P_L[(\p^\mu\b{Z})\p_\mu\chi + \gamma^{\mu\nu}\p_\nu(\p_\mu\b{Z}\chi)]
\end{equation}
in \eqref{eq: variazione chirale left azione cinetica}. Due dei quattro termini si cancellano. Dopo l'integrazione per parti e che $\eta^{\mu\nu} - \gamma^{\mu\nu} = \gamma^\nu\gamma^\mu$, si ottiene 
\begin{equation}
\label{eq: variazione azione cinetica (finale)}
    \delta_{P_L\epsilon}S_K = -\frac{1}{\sqrt{2}}\int d^4x\,\p_\mu\b{\epsilon}P_L(\sl{\p}\b{Z})\gamma^\mu\chi.
\end{equation}
Questo dimostra che $\delta_{P_L\epsilon}S_K$ si annulla per uno spinore costante $\epsilon$, verificando così che l'azione è supersimmetrica. 

Passando al termine di interazione $S_F$, la variazione $\delta S_F$ sotto le trasformazioni in \eqref{eq: trasformazioni multipletto chirale} ha la struttura 
\begin{equation}
    \begin{split}
        \delta S_F &= \int d^4x\,\bigg[\delta FW'(Z) + \delta ZFW''(Z) \\
        &- \delta\b{\chi}P_L\chi W''(Z) - \frac{1}{2}\delta Z\b{\chi}P_L\chi W'''(Z)\bigg],
    \end{split}
\end{equation}
dove si sono considerate le derivate di $W(Z)$ richieste per includere tutte le sorgenti della variazione $\delta Z$. Dopo aver usato \eqref{eq: trasformazioni multipletto chirale} si combinano i termini. I due termini in $FP_L\chi$ si cancellano e si rimane con 
\begin{equation}
\label{eq: variazione azione interagente (finale)}
    \delta S_F = \frac{1}{\sqrt{2}}\int d^4x\,\bigg[\b{\epsilon}\sl{\p}(W'P_L\chi) - \frac{1}{2}W'''\b{\epsilon}P_L\chi\b{\chi}P_L\chi\bigg],
\end{equation}
dove l'ultimo termine si annulla per le proprietà di anti-commutazione dato che è il prodotto di tre spinori. Quindi, anche $\delta S_F$ si annulla per uno spinore $\epsilon$ costante ed è supersimmetrica. Chiaramente, si annulla anche $\delta S_{\b{F}}$ essendo semplicemente il coniugato di $\delta S_F$. Dunque, la supersimmetria è stabilita al livello interagente. 

I termini in $\p_\mu\b{\epsilon}$ rimanenti in \eqref{eq: variazione azione cinetica (finale)} e \eqref{eq: variazione azione interagente (finale)} insiemi ai loro coniugati si combinano per formare la supercorrente di Noether della teoria interagente. Questa è 
\begin{equation}
    J^\mu = \frac{1}{\sqrt{2}}[P_L(\sl{\p}\b{Z} - F) + P_R(\sl{\p}Z - \b{F})]\gamma^\mu\chi,
\end{equation}
dove si sono usate le equazioni del moto per il campo ausiliario $F = -\b{W}'(\b{Z})$ e $\b{F} = -W'(Z)$.

\section{CONTO A LEZIONE}
La variazione totale dell'azione $S_K$ la si può scrivere come 
\begin{equation}
    \begin{split}
        \sqrt{2}\delta S_K &= \sqrt{2}\delta_{P_L\epsilon} S_K + \sqrt{2}\delta_{P_R\epsilon} S_K \\
        &= \int d^4x\,[-\p^\mu \b{Z}\p_\mu(\epsilon P_L\chi) + \b{\epsilon}P_L\sl{\p}\b{Z}\sl{\p}P_L\chi \\
        &- \b{\chi}P_R\sl{\p}(P_LF\epsilon) + F\b{\epsilon}P_L\sl{\p}\chi].
    \end{split} 
\end{equation}
Adesso, riscrivendo l'ultimo termine come
\begin{equation}
    \b{\epsilon}P_L\sl{\p}\chi = \epsilon^TCP_L\gamma^\mu\p_\mu\chi = \overline{P_L\epsilon}\gamma^\mu\p_\mu\chi = -\p_\mu\b{\chi}\gamma^\mu P_L\epsilon
\end{equation}
e combinandolo con $\b{\chi}P_R\sl{\p}(P_LF\epsilon)$, si ottiene una derivata totale
\begin{equation}
    -\p_\mu(\b{\chi}\gamma^\mu P_LF\epsilon),
\end{equation}
che sostituita in $\delta S_K$ porta a
\begin{equation}
    \delta S_K = \int d^4x\,[-\p^\mu\b{Z}(\p_\mu \b{\epsilon}P_L\chi + \b{\epsilon}P_L\p_\mu\chi) + \b{\epsilon}P_L\sl{\p}\b{Z}\sl{\p}P_L\chi].
\end{equation}
Sfruttando ora l'identità
\begin{equation}
    \b{\epsilon}P_L\sl{\p}\b{Z}\sl{\p}\chi = \b{\epsilon}P_L\gamma^\mu\gamma^\nu\p_\nu\chi\p_\mu\b{Z} = \b{\epsilon}P_L\p^\mu + \b{\epsilon}P_L\gamma^{\mu\nu}\p_\nu(\p_\mu \b{Z}\chi),
\end{equation}
dove si è tenuto conto che $\gamma^{\mu\nu}\p_\mu\p_\nu Z = 0$ e si è usato che $\gamma^\mu\gamma^\nu = \eta^{\mu\nu} + \gamma^{\mu\nu}$, si ha 
\begin{equation}
    \delta S_K = \int d^4x\,[-\p^\mu\b{Z}(\p_\mu\b{\epsilon}P_L\chi) + \b{\epsilon}P_L\gamma^{\mu\nu}\p_\nu(\p_\mu\b{Z}\chi)].
\end{equation}
Integrando adesso per parti il secondo termine, si ottiene
\begin{equation}
    \begin{split}
        \delta S_K &= \int d^4x\,[-\p_\mu\b{Z}(\p_\mu\b{\epsilon}P_L\chi) - (\p_\nu\b{\epsilon})P_L\gamma^{\mu\nu}\p_\mu\b{Z}\chi] \\
        &= \int d^4x\, [\p_\nu\b{\epsilon}P_L\gamma^\nu\gamma^\mu\chi\p_\mu\b{Z}], 
    \end{split}
\end{equation}
che effettivamente soddisfa $\delta S_K = 0$ quando $\epsilon$ è costante. Per determinare la corrente conservata, si deve temporaneamente supporre che il parametro $\epsilon$ non sia costante, in questo modo, dopo la variazione, si trova che 
\begin{equation}
    J^\mu_K = \frac{1}{\sqrt{2}}(P_L\sl{\p}\b{Z}\gamma^\mu\chi + P_R\sl{\p}Z\gamma^\mu\chi),
\end{equation}
dove 
\begin{equation}
    (\p_\mu\b{\epsilon}P_L\sl{\p}\b{Z}\gamma^\mu\chi)^* = (\p_\mu\b{\epsilon}P_L\sl{\p}\b{Z}\gamma^\mu\chi)^C = (\p_\mu\b{\epsilon})^CP_L^C\sl{\p}^CZ(\gamma^{\mu})^C\chi.
\end{equation}
Si noti che la corrente $J^\mu_K$ ha anche struttura di spinore data da $\chi$, quindi è uno spinore-vettore. Dalla legge di conservazione $\p_\mu J_K^\mu = 0$ si trova che la carica conservata $Q_K$, detta \emph{supercarica}, è uno spinore in quanto si è fissato un indice vettoriale di Lorentz. 

L'azione completa del termine di interazione è $S = S_K + S_F + S_{\b{F}}$, dove 
\begin{equation}
    S_F = \int d^4x\,FW'(Z) - \frac{1}{2}\b{\chi}P_L\chi W''(Z),
\end{equation}
con $W(Z)$ una funzione olomorfa che rappresenta un \emph{superpotenziale} al momento arbitrario. Le variazioni sono 
\begin{equation}
    \delta \b{\chi}P_L = \frac{1}{\sqrt{2}}\b{\epsilon}(-\sl{\p}Z + F)P_L,\quad \delta W' = W''\delta Z,\quad \delta W'' = W'''\delta Z.
\end{equation}
Facendo la variazione tutti i termini in $F$ si cancellano, rimane 
\begin{equation}
    \begin{split}
        S_F &= \int d^4x\,\bigg[\b{\epsilon}\sl{\p}P_L\chi W'(Z) - \frac{1}{2}W'''(Z)\b{\chi}P_L\b{\epsilon}P_L\chi \\
        &+ \frac{1}{2}W''(Z)(\b{\epsilon}P_R\sl{\p}Z\chi - \b{\chi}P_L\sl{\p}Z\epsilon)\bigg],
    \end{split}
\end{equation}
dove il secondo termine si annulla in $D = 4$ dato che contiene tre spinori $\chi$. Effettuando il Majorana flip $-\b{\chi}P_L\sl{\p}Z\epsilon = \b{\epsilon}P_R\sl{\p}Z\chi$, dove si è usato $\b{\chi}P_L = \bigb{P_R\chi}$, si calcola 
\begin{equation}
    \p_\mu(\b{\epsilon}\gamma^\mu W'P_L\chi) = \p_\mu\b{\epsilon}\gamma^\mu W'P_L\chi + \b{\epsilon}\sl{\p} ZW''P_L\chi + \b{\epsilon}\gamma^\mu W'P_L\p_\mu\chi.
\end{equation}
Allora la somma dei due è una derivata totale, che porta ad un integrale nullo, più il termine 
\begin{equation}
    -\int d^4x\,\p_\mu\b{\epsilon}\gamma^\mu W'P_L\chi = 0,
\end{equation}
che si annulla pure dato che $\epsilon$ è assunto costante. Dunque, la supercorrente di Noether è 
\begin{equation}
    J^\mu = J^\mu_K + [\gamma^\mu W'P_L\chi + (\gamma^\mu W'P_L\chi)^*].
\end{equation}